               1. Bank schema:
Branch(BranchName, BranchCity, Assets) Customer(CustomerID,
CustomerName, CustomerStreet, CustomerCity)
Account(AccountNumber, BranchName, Balance) 
Loan(LoanNumber,BranchName, Amount)
Depositor(CustomerID, AccountNumber)
Borrower(CustomerID, LoanNumber)


1. Retrieve distinct customers who have either an account or a loan:
SELECT DISTINCT CustomerID, CustomerName
FROM Customer
WHERE CustomerID IN (
SELECT CustomerID FROM Depositor
UNION
SELECT CustomerID FROM Borrower
);

2. Retrieve the branch with the maximum number of accounts:
SELECT BranchName
FROM Account
GROUP BY BranchName
HAVING COUNT(*) = (
SELECT MAX(AccountCount)
FROM (
SELECT COUNT(*) AS AccountCount
FROM Account
GROUP BY BranchName
) AS BranchCounts
);

3. Find the maximum loan amount issued in each branch:
SELECT BranchName, MAX(Amount) AS MaxLoan
FROM Loan
GROUP BY BranchName;

4. Find customers whose name starts with 'A'
SELECT CustomerID, CustomerName
FROM Customer
WHERE CustomerName LIKE 'A%';

5. Delete all loan amounts between 5000 and 15000:
DELETE FROM Loan
WHERE Amount BETWEEN 5000 AND 15000;

6. Find branches that have no accounts:
SELECT BranchName
FROM Branch
WHERE BranchName NOT IN (
SELECT DISTINCT BranchName FROM Account
);

7. Find branches where total loan amount is greater than 50,000:
SELECT BranchName, SUM(Amount) AS TotalLoan
FROM Loan
GROUP BY BranchName
HAVING SUM(Amount) > 50000;

8. Deduct service charge: 3% if customer has both loan and account, else 5%:
UPDATE Account
SET Balance = Balance * (
CASE
WHEN AccountNumber IN (
SELECT d.AccountNumber
FROM Depositor d
JOIN Borrower b ON d.CustomerID = b.CustomerID
) THEN 0.97
ELSE 0.95
END
);
